\section{Descent and stability}\label{sec:descent}

We iterate the subdivision estimate while the untilted growth or the
tilt is large. The first goal is a bound, independent of the number of
levels, for the total residual growth and the accumulated changes of
tilt. This bound will show that a sufficiently large initial tilt with
small residual growth forces high untilted growth to persist on a
positive fraction of paths.

Fix the integer \(A\) and constants \(C_0,C_1\) from
Proposition~\ref{prop:subdivision}.  We may increase \(C_0\) so that
\(C_0\ge 1\).
Consider a finite subdivision tree with all terminal squares at the same
level, and suppose that
\[
K_Q=k s(Q)\ge a_0>0
\]
on every square of the tree.  The square at the start of a subtree will
be denoted by \(I\), and its side satisfies \(s(I)\le 1\).
All disks used below are assumed to lie in the coordinate region of
Section~\ref{sec:subdivision}.  The arguments also apply to a subtree
in an extension of the same grid whenever its disks have this property.
Throughout this section, constants may depend on the fixed number
\(a_0\), as well as on the coordinate data and the constants already
chosen.

Choosing a child uniformly among the \(A^2\) children at each step
defines a random path through the tree.  Equivalently, one may choose a
point uniformly in \(I\) and follow its containing squares, ignoring
grid boundaries of area zero.

\begin{lemma}[One step of descent]\label{lem:descent}
There is a fixed \(L>0\) such that, at any nonterminal square \(Q\),
if \(N(Q)>L\) or \(\abs{b}>L\), one can assign a tilt \(b_q\) to each
child \(q\) so that, writing
\(m=m(Q,b)\), \(m'=m(q,b_q)\), and \(b'=b_q\), one has
\begin{equation}\label{eq:descent}
 \begin{gathered}
 \E\bigl(\abs{b'-b}\mid Q,b\bigr)\le C_1(m+f_Q),
 \qquad
 \E\bigl(m'\mid Q,b\bigr)\le \tfrac12(m+f_Q),\\
 f_Q=\frac{C_0}{K_Q}+s(Q).
 \end{gathered}
\end{equation}
Moreover, if \(Q_0=I,Q_1,\ldots,Q_d\) is any path from \(I\) to
the terminal level, then
\begin{equation}\label{eq:floors}
 \sum_{i=0}^{d}f_{Q_i}
 \le \frac{A}{A-1}
       \left(\frac{C_0}{a_0}+s(I)\right).
\end{equation}
\end{lemma}

\begin{proof}
Write \(K=K_Q\), \(s=s(Q)\), and set
\[
 S=E(Q,b)+C_0+Ks,\qquad B=Kb.
\]
Thus \(S\ge C_0\), \(E(Q,b)\le S\), and \(S\ge Ks\).
Equation~\eqref{eq:untilted-comparison} gives
\[
 N(Q)\le m(Q,b)+(R+1)\abs{b}
       \le (R+1)\bigl(m(Q,b)+\abs{b}\bigr).
\]
Consequently \(S+\abs{B}\ge KN(Q)/(R+1)\); also
\(S+\abs{B}\ge K\abs{b}\).
Choose, for example, \(L=2C_0(R+1)\).
Either of the alternatives in the lemma then implies
\[
 \frac{K}{S+\abs{B}}\le \frac1{C_0}.
\]
Using \(S\ge Ks\), we obtain the other required smallness condition,
\[
 \frac{K^2s}{S(S+\abs{B})}
 \le \frac{K}{S+\abs{B}}\le \frac1{C_0}.
\]
All the hypotheses of Proposition~\ref{prop:subdivision} hold.
Its conclusions, with \(S/K=m(Q,b)+f_Q\), are precisely
\eqref{eq:descent}.

Along a path, \(s(Q_i)=A^{-i}s(I)\) and
\(K_{Q_i}=A^{d-i}K_{Q_d}\).  Since \(K_{Q_d}\ge a_0\),
\[
 \sum_{i=0}^d \frac{C_0}{K_{Q_i}}
 \le \frac{C_0}{a_0}\sum_{\ell=0}^d A^{-\ell},
 \qquad
 \sum_{i=0}^d s(Q_i)
 =s(I)\sum_{i=0}^d A^{-i}.
\]
The geometric-series bound gives \eqref{eq:floors}.
\end{proof}

We next allow descent to stop according to the values encountered
along a path.  A stopping rule is understood to use only the squares
and tilts already visited.  In particular, a violation first detected
on arrival at a child stops the procedure at that child.

\begin{lemma}[Stopped descent]\label{lem:stopped-descent}
Start at \(I\) with a prescribed tilt \(b_I\).  Whenever a step is
made, choose the child tilts by Lemma~\ref{lem:descent}, and suppose
that its applicability condition holds at every square from which
the procedure continues.  Stop at a time \(T\), no later than the
terminal level.  Write \(m_i=m(Q_i,b_i)\) for the visited values,
including the last value \(m_T\).  There is a constant
\(C_{\mathrm{st}}\), independent of the depth and of the stopping
rule, such that
\begin{equation}\label{eq:stopped-descent}
 \E\left[
      \sum_{i=0}^{T}m_i
      +\sum_{i=0}^{T-1}\abs{b_{i+1}-b_i}
    \right]
 \le C_{\mathrm{st}}\bigl(m(I,b_I)+1\bigr).
\end{equation}
The second sum includes the step producing the last square, even if
that square causes the procedure to stop.
\end{lemma}

\begin{proof}
There are only finitely many nodes.  At each node from which descent
continues, choose one of the assignments furnished by
Lemma~\ref{lem:descent}.  These choices and the stopping rule determine
all random variables on the finite probability space of terminal
paths.  Thus only finite conditional-expectation identities are
needed.  In expressions with a fixed index, set \(m_i=0\) after
stopping; those auxiliary values do not enter the stopped sums.

Let \(d\) be the depth below \(I\), put \(m_0=m(I,b_I)\), and set
\[
 F_I=\frac{A}{A-1}\left(\frac{C_0}{a_0}+s(I)\right),
 \qquad
 M=\E\sum_{i=0}^{T}m_i.
\]
The event \(\{i<T\}\) is determined before the next child is selected.
Applying \eqref{eq:descent} on this event and summing over
\(0\le i<d\) gives
\[
 M-m_0
 =\sum_{i=0}^{d-1}\E\bigl[\mathbf 1_{\{i<T\}}m_{i+1}\bigr]
 \le \frac12\E\sum_{i=0}^{T-1}(m_i+f_{Q_i})
 \le \frac12 M+\frac12 F_I.
\]
All the \(m_i\) are nonnegative by Definition~\ref{def:growth}.
It follows that \(M\le 2m_0+F_I\).
The tilt estimate in \eqref{eq:descent} similarly yields
\[
 \E\sum_{i=0}^{T-1}\abs{b_{i+1}-b_i}
 \le C_1\E\sum_{i=0}^{T-1}(m_i+f_{Q_i})
 \le 2C_1m_0+2C_1F_I.
\]
Their sum is at most
\[
 2(1+C_1)m_0+(1+2C_1)F_I.
\]
Since \(s(I)\le1\), this is bounded by the right-hand side of
\eqref{eq:stopped-descent} with a constant depending only on the fixed
data.  This calculation includes \(T=0\), and it includes both \(m_T\)
and the incoming increment when the last square violates a stopping
condition.
\end{proof}

A large tilt whose remaining excess is small forces the original,
untilted growth to be large.  We give a separate size threshold and
ratio threshold for this statement.

\begin{lemma}[Large slope forces large growth]\label{lem:large-slope}
For the fixed \(L\) in Lemma~\ref{lem:descent}, there are constants
\(B_*\ge1\) and \(\eta_*>0\) such that every admissible square with
\(s(Q)\le1\) and \(K_Q\ge a_0\) satisfies
\begin{equation}\label{eq:large-slope}
 \abs{b}\ge B_*,
 \qquad m(Q,b)\le\eta_*\abs{b}
 \quad\Longrightarrow\quad N(Q)>L.
\end{equation}
\end{lemma}

\begin{proof}
If no such pair of thresholds existed, then for each positive integer
\(j\) there would be a square \(Q_j\) and a tilt \(b_j\) with
\[
 \abs{b_j}\ge j,\qquad
 \frac{m(Q_j,b_j)}{\abs{b_j}}\le\frac1j,\qquad
 N(Q_j)\le L.
\]
Rescale about the center of \(Q_j\) by its side \(s_j=s(Q_j)\).
In the rescaled coordinates, \(U_j\) solves
\(\Delta U_j+K_j^2p_jU_j=0\) on \(\D_R\), where
\(K_j=K_{Q_j}\), \(\norm{p_j}_\infty\le C\), and
\(\norm{\nabla p_j}_\infty\le Cs_j\).
Set
\[
 S_j=K_j\abs{b_j},\qquad B_j=K_jb_j,\qquad
 D_j=S_j+\abs{B_j}=2S_j.
\]
The hypotheses of Theorem~\ref{thm:subharmonic-profile} follow from
\[
 S_j\ge a_0j\longrightarrow\infty,\qquad
 \frac{K_j}{D_j}=\frac1{2\abs{b_j}}\longrightarrow0,\qquad
 \frac{K_j^2\norm{\nabla p_j}_\infty}{S_jD_j}
 \le\frac{Cs_j}{2\abs{b_j}^2}\longrightarrow0.
\]
Normalize the tilted measures
\[
 d\mu_j(x)=c_j e^{-2B_j\cdot x}\abs{U_j(x)}^2\,dx
\]
to have mass one on \(\D_R\).  By
Lemma~\ref{lem:profile-extraction} and
Theorem~\ref{thm:subharmonic-profile}, a subsequence has a
subharmonic profile \(V\le0\), with the identically minus infinity
case initially allowed.  The definitions give
\[
 \frac{\log\mu_j(\D_1)}{2S_j}
 =-\frac{E(Q_j,b_j)}{S_j}
 =-\frac{m(Q_j,b_j)}{\abs{b_j}}\longrightarrow0.
\]
The upper bound in \eqref{eq:profile-bounds} on the compact set
\(\overline{\D_1}\) therefore implies
\(\sup_{\overline{\D_1}}V=0\).
In particular the profile is nontrivial and attains its maximum at
an interior point of \(\D_R\).  The strong maximum principle yields
\(V=0\) on \(\D_R\).

Pass to a further subsequence so that
\(e_j=b_j/\abs{b_j}\) converges to a unit vector \(e\).
Undoing the tilt expresses the untilted growth as
\[
 \frac{E(Q_j,0)}{S_j}
 =
 \frac{1}{2S_j}\log
       \int_{\D_R}e^{2S_je_j\cdot x}\,d\mu_j(x)
 -
 \frac{1}{2S_j}\log
       \int_{\D_1}e^{2S_je_j\cdot x}\,d\mu_j(x).
\]
For any fixed \(1<r<R\), bound the first integral from below by its
restriction to \(\D_r\).  The weighted form of
Lemma~\ref{lem:profile-extraction}, applied from below on the open
disk \(\D_r\) and from above on the compact disk
\(\overline{\D_1}\), gives
\[
 \liminf_{j\to\infty}\frac{E(Q_j,0)}{S_j}
 \ge
 \sup_{\D_r} e\cdot x
 -\sup_{\overline{\D_1}} e\cdot x
 =r-1.
\]
Letting \(r\uparrow R\) gives a lower bound of \(R-1\), without any
claim about mass at the boundary of \(\D_R\).  On the other hand,
\[
 0\le\frac{E(Q_j,0)}{S_j}
 =\frac{N(Q_j)}{\abs{b_j}}
 \le\frac{L}{j}\longrightarrow0,
\]
a contradiction.
\end{proof}

\begin{proposition}[Stability of large growth]\label{prop:stability}
There are constants \(P\ge1\) and \(\sigma>0\) with the following
property.  Suppose a square \(I\) and a tilt \(b_I\) satisfy
\begin{equation}\label{eq:stability-hypotheses}
 \abs{b_I}\ge P,\qquad m(I,b_I)\le\sigma\abs{b_I}.
\end{equation}
Then at least half the paths from \(I\) to the terminal level satisfy
\(N(Q)>L\) at every visited square \(Q\), including \(I\) and the
terminal square.  Equivalently, these paths occupy a measurable set
of area at least \(\abs{I}/2\).
\end{proposition}

\begin{proof}
Let \(B_*,\eta_*\) be the thresholds from
Lemma~\ref{lem:large-slope}.  Choose \(\gamma\) so that
\[
 0<\gamma<\frac12,\qquad
 \frac{\gamma}{1-\gamma}\le\eta_*.
\]
Let \(C_{\mathrm{st}}\) be the constant from
Lemma~\ref{lem:stopped-descent}.  Choose
\[
 0<\sigma\le\min\left\{\frac{\gamma}{2},
                         \frac{\gamma}{8C_{\mathrm{st}}}\right\},
 \qquad
 P\ge\max\left\{1,\frac{B_*}{1-\gamma},
                  \frac{2L}{1-\gamma},
                  \frac{8C_{\mathrm{st}}}{\gamma}\right\}.
\]
Write \(B_0=\abs{b_I}\).  Starting at \(I\), descend using
Lemma~\ref{lem:descent} while
\[
 \abs{b_i-b_I}<\gamma B_0,\qquad m_i<\gamma B_0,
\]
and stop at the first failure of either inequality or at the terminal
level, whichever comes first.  These inequalities hold initially by
\eqref{eq:stability-hypotheses}.  Whenever they hold, we have
\[
 \abs{b_i}>(1-\gamma)B_0\ge B_*,
 \qquad \abs{b_i}>L,
 \qquad
 \frac{m_i}{\abs{b_i}}
 <\frac{\gamma}{1-\gamma}\le\eta_*.
\]
Thus every required step is available and
Lemma~\ref{lem:large-slope} gives \(N(Q_i)>L\) at each such square.

Let \(T\) be this stopping time.  Include in the failure event a
violation first produced at the terminal level.
If a failure occurs, then either \(m_T\ge\gamma B_0\), or the
triangle inequality gives
\(\sum_{i=0}^{T-1}\abs{b_{i+1}-b_i}\ge\gamma B_0\).
Consequently Lemma~\ref{lem:stopped-descent} and Markov's inequality
give
\[
 \begin{split}
 \mathbb P(\text{failure})
 &\le
 \frac{1}{\gamma B_0}
 \E\left[\sum_{i=0}^{T}m_i+
         \sum_{i=0}^{T-1}\abs{b_{i+1}-b_i}\right]\\
 &\le\frac{C_{\mathrm{st}}}{\gamma}
                \left(\sigma+\frac1P\right)
 \le\frac14.
 \end{split}
\]
On the complementary event, descent reaches the terminal level and
the two strict inequalities hold at every visited square.
All of those squares have \(N>L\).  This event has probability at
least \(3/4\), which proves the asserted lower bound of \(1/2\).
The same reasoning includes a subtree of depth zero.
\end{proof}
